% =====================================================================
%  FICHE 11 — THEME 5 — Exercices FACILE — Piles, Nernst, Faraday
% =====================================================================
\documentclass[11pt,a4paper]{article}
\usepackage[utf8]{inputenc}
\usepackage[T1]{fontenc}
\usepackage{lmodern}
\usepackage[french]{babel}
\usepackage{amsmath,amssymb}
\usepackage[version=4]{mhchem}
\usepackage{siunitx}
\sisetup{output-decimal-marker={,},per-mode=symbol}
\usepackage{xcolor}
\usepackage{array}
\usepackage{booktabs}
\usepackage{enumitem}
\usepackage[most]{tcolorbox}
\usepackage[a4paper,margin=2.1cm,headheight=26pt,headsep=10pt]{geometry}
\usepackage{fancyhdr}
\usepackage[hidelinks]{hyperref}

\definecolor{primary}{RGB}{146,95,160}
\definecolor{primarydark}{RGB}{92,52,104}
\newcommand{\NO}{\ensuremath{\mathrm{N.O.}}}
\newcommand{\Ered}{\ensuremath{E_{\mathrm{r}}^{\circ}}}
\newcommand{\fem}{\ensuremath{\Delta E}}

\newcounter{exo}
\newtcolorbox{exo}[1][]{enhanced,breakable,boxrule=0.9pt,arc=2pt,colback=primary!4,
  colframe=primary,fonttitle=\bfseries,coltitle=white,left=3mm,right=3mm,top=2.2mm,bottom=2.2mm,
  title={Exercice~\refstepcounter{exo}\theexo\ifx\relax#1\relax\relax\else\ \textmd{--- \itshape #1}\fi}}
\newtcolorbox{donnees}{enhanced,boxrule=0.7pt,arc=2pt,colback=black!3,colframe=black!45,
  left=3mm,right=3mm,top=2mm,bottom=2mm,fonttitle=\bfseries\small,coltitle=black!70,title={Données}}

\pagestyle{fancy}\fancyhf{}
\fancyhead[L]{\small\textcolor{primarydark}{\textbf{Fiche 11 · Thème 5} — Piles, Nernst, Faraday}}
\fancyhead[R]{\small\textcolor{primarydark}{\textsc{Facile}}}
\fancyfoot[C]{\small\thepage}
\renewcommand{\headrulewidth}{0.8pt}
\renewcommand{\headrule}{\hbox to\headwidth{\color{primary}\leaders\hrule height \headrulewidth\hfill}}

\begin{document}
\begin{center}
{\LARGE\bfseries\color{primarydark}Thème 5 — Niveau Facile}\\[2pt]
{\small\itshape Anode/cathode, notation, f.é.m. standard et première loi de Faraday.}\\[-4pt]
{\color{primary}\rule{\textwidth}{1pt}}
\end{center}

\begin{donnees}
\centering\small
$\Ered(\ce{Ag+}/\ce{Ag})=+0{,}80$\,V ; \quad $\Ered(\ce{Cu^{2+}}/\ce{Cu})=+0{,}34$\,V ; \quad
$\Ered(\ce{Zn^{2+}}/\ce{Zn})=-0{,}76$\,V ; \quad $F=\SI{96485}{C\per mol}$.
\end{donnees}

\begin{exo}[anode ou cathode ?]
On réalise la pile associant les couples \ce{Cu^{2+}}/\ce{Cu} et \ce{Zn^{2+}}/\ce{Zn}.
\begin{enumerate}[label=(\alph*),leftmargin=2em,topsep=1pt,itemsep=1pt]
  \item Quelle électrode est l'\textbf{anode} ($-$) ? Laquelle est la \textbf{cathode} ($+$) ?
  \item Écris la demi-équation qui se produit à chaque électrode (oxydation / réduction).
\end{enumerate}
\end{exo}

\begin{exo}[circulation des porteurs de charge]
Pour la pile \ce{Zn}$\,\|\,$\ce{Cu} qui débite :
\begin{enumerate}[label=(\alph*),leftmargin=2em,topsep=1pt,itemsep=1pt]
  \item Dans le fil extérieur, dans quel sens circulent les \textbf{électrons} (de l'anode vers la
        cathode, ou l'inverse) ?
  \item À quoi sert le \textbf{pont salin} ? Dans quel sens migrent les anions ?
\end{enumerate}
\end{exo}

\begin{exo}[écrire la notation conventionnelle]
On construit la pile où le zinc est l'anode et l'argent la cathode, avec les ions \ce{Zn^{2+}} et
\ce{Ag+} en solution. Écris la \textbf{représentation conventionnelle} de cette pile (avec les
symboles \textbar\ et $\|$ et les pôles $-$ et $+$).
\end{exo}

\begin{exo}[force électromotrice standard]
Calcule la f.é.m. standard $\fem^{\circ}$ de la pile \ce{Zn}$\,\|\,$\ce{Ag} (conditions standard),
sachant que la cathode est l'argent et l'anode le zinc. Utilise
$\fem^{\circ}=\Ered(\text{cathode})-\Ered(\text{anode})$.
\end{exo}

\begin{exo}[première approche de la loi de Faraday]
Un courant d'intensité $I=\SI{2.0}{A}$ traverse une cuve d'électrolyse pendant $t=\SI{1930}{s}$.
\begin{enumerate}[label=(\alph*),leftmargin=2em,topsep=1pt,itemsep=1pt]
  \item Calcule la charge électrique totale $Q=It$ (en coulombs).
  \item Combien de \textbf{moles d'électrons} cela représente-t-il ($n_{\ce{e^-}}=Q/F$) ?
\end{enumerate}
\end{exo}

\end{document}
